SimpleChef is a mobile-first cooking assistant designed to reduce friction across the entire cooking process---from discovering a recipe to meal planning and grocery shopping. The tagline \emph{Your calm kitchen companion} reflects the goal of a capable, calm tool that streamlines everyday cooking tasks.

As stated in our project proposal, many cooking apps fragment the experience: recipes in one place, calories in another, grocery lists on paper or in notes. Instructions often appear as long blocks of text, and timers frequently push users out of the recipe context. SimpleChef targets a single cohesive flow: recipe discovery, step-by-step cooking with integrated timers, meal planning with calorie awareness, and grocery support. The design applies HCI ideas including progressive disclosure, clarity over clutter, and smart defaults with full user control~\cite{preece2015,hartson2012ux}.

\subsection{Objectives (aligned with the proposal)}
\begin{itemize}[leftmargin=*]
  \item \textbf{Unify the workflow} across browsing, cooking mode, planning, and groceries with consistent navigation and data flow.
  \item \textbf{Optimize step-by-step cooking} so each step is scannable during active cooking, with integrated timers and (in the full design) mise en place tied to steps.
  \item \textbf{Timer management} via a timer dock that surfaces active timers without leaving the cooking context.
  \item \textbf{Grocery automation} from planned meals (deduplication, categories, editing, export) as a later integration milestone.
  \item \textbf{Lower recipe input friction} via manual entry plus AI-assisted text parsing with verification before save; image and video import remain planned.
  \item \textbf{Meal planning and calorie goals} by assigning meals to days and relating plans to profile goals; dietary filtering and richer statistics are extensions.
  \item \textbf{Iterate with users} through feedback sessions and refinement of flows and visuals~\cite{nielsen1994heuristic}.
\end{itemize}

\subsection{Tools and technology}
\begin{itemize}[leftmargin=*]
  \item \textbf{Frontend:} React Native (Expo) for cross-platform mobile development. Primary screens live under \texttt{frontend/app/} (e.g., tab routes for Home, Calendar, Add, Grocery, Profile; stack routes for recipe detail, cooking mode, manual entry).
  \item \textbf{Backend:} FastAPI (Python) with SQLAlchemy models and Alembic migrations; JWT auth for protected routes (\texttt{backend/app/api/}).
  \item \textbf{Data store:} PostgreSQL (via Docker Compose for local development).
  \item \textbf{Design:} Figma for wireframes and high-fidelity exploration~\cite{figma2026help}.
  \item \textbf{AI/ML:} Planned for parsing text, images, or video; the current server uses a structured placeholder parser in \texttt{backend/app/services/ai\_service.py} to exercise parse-then-edit UX end to end.
\end{itemize}
